\chapter{Estado da Arte}\label{cap:1-est-art}

O avanço das tecnologias digitais e a crescente preocupação com a saúde e o bem-estar têm impulsionado o desenvolvimento de sistemas inteligentes para análise de dados nutricionais. O estado da arte nesta área abrange a interseção entre sistemas de informação em saúde, inteligência artificial, aprendizado de máquina, processamento de linguagem natural e a integração de modelos de linguagem de larga escala (LLM) para suporte à tomada de decisão nutricional personalizada.

\section{Sistemas de Informação em Saúde e Nutrição}

Sistemas de informação em saúde têm sido amplamente estudados como ferramentas para apoiar o monitoramento, a análise e a gestão de dados clínicos e nutricionais \cite{hersh2022health}. Tais sistemas buscam não apenas registrar informações, mas também fornecer análises preditivas e recomendações personalizadas, promovendo a prevenção de doenças e a melhoria da qualidade de vida. No contexto nutricional, plataformas digitais vêm sendo utilizadas para registro de refeições, cálculo de ingestão calórica, acompanhamento de dietas e geração de relatórios para profissionais e pacientes \cite{martins2021aplicativos}.

\section{Inteligência Artificial e Modelos de Linguagem em Nutrição}

A aplicação de inteligência artificial (IA) em saúde e nutrição tem se destacado nos últimos anos, especialmente com o uso de aprendizado de máquina para análise de grandes volumes de dados e identificação de padrões alimentares \cite{topol2019deep, chen2020artificial}. Recentemente, modelos de linguagem de larga escala (LLM), como o GPT-4 \cite{openai2023gpt4}, têm sido explorados para tarefas de processamento de linguagem natural, extração de informações de textos clínicos e geração de recomendações automáticas \cite{singhal2023large}. A integração desses modelos em sistemas de análise nutricional permite consultas inteligentes, respostas em linguagem natural e maior personalização no atendimento ao usuário \cite{wang2023llmhealth}.

\section{Rastreamento, Verificação e Validação de Requisitos em Sistemas de Saúde}

A engenharia de requisitos para sistemas de saúde demanda rigor em rastreabilidade, verificação e validação, devido à criticidade dos dados e à necessidade de conformidade com normas e padrões \cite{gotel1994traceability}. Estudos apontam que a adoção de matrizes de rastreabilidade, técnicas de verificação sistemática e validação com stakeholders são essenciais para garantir a qualidade, a segurança e a aceitação dos sistemas \cite{goncalves2004ieee830, pinheiro2003requirements}. A literatura também destaca a importância de requisitos não funcionais, como desempenho, segurança, usabilidade e escalabilidade, especialmente em ambientes que lidam com dados sensíveis e integração com serviços externos \cite{sommerville2003engenharia}.

\section{Tendências e Desafios Atuais}

Entre as tendências atuais, destaca-se a adoção de arquiteturas baseadas em microserviços, uso de APIs para integração com bancos de dados alimentares e serviços de IA, além da preocupação crescente com privacidade e proteção de dados \cite{deng2021microservices, fernandes2022privacidade}. Desafios incluem a interoperabilidade entre sistemas, a explicabilidade dos modelos de IA, a adaptação a diferentes perfis de usuários e a validação clínica das recomendações geradas \cite{holzinger2017we, doshi2017towards}.

\section{Considerações Finais}

O estado da arte evidencia que sistemas de análise de dados nutricionais integrados a LLM representam uma fronteira inovadora, com potencial para transformar a prática nutricional e o autocuidado. No entanto, exigem abordagens robustas de engenharia de requisitos, validação contínua e atenção às questões éticas e regulatórias.

% \textbf{Palavras-chave:} sistemas de informação em saúde, análise de dados nutricionais, inteligência artificial, modelos de linguagem, requisitos de software, rastreabilidade, validação, privacidade, usabilidade, microserviços.
